% env-doc.tex - Environment-Definitionen für die Dokumentation
% ConTeXt MkIV / LuaMetaTeX

% =============================================================================
% CODE-BLÖCKE (für Python, Bash, etc.)
% =============================================================================

\definetyping[pythoncode][
    before={\blank[small]\startframedtext[
        frame=on,
        framecolor=codeborder,
        background=color,
        backgroundcolor=codebackground,
        width=\textwidth,
        align=flushleft,
    ]},
    after={\stopframedtext\blank[small]},
    style=\tt\small,
]

\definetyping[bashcode][
    before={\blank[small]\startframedtext[
        frame=on,
        framecolor=codeborder,
        background=color,
        backgroundcolor=codebackground,
        width=\textwidth,
        align=flushleft,
    ]},
    after={\stopframedtext\blank[small]},
    style=\tt\small,
]

% Inline-Code
\definetype[code][style=\tt\small]

% =============================================================================
% HINWEIS-BOXEN
% =============================================================================

% Info-Box (blau)
\defineframedtext[infobox][
    frame=on,
    framecolor=primarycolor,
    background=color,
    backgroundcolor=primarycolor:1,
    width=\textwidth,
    before={\blank[medium]},
    after={\blank[medium]},
    loffset=1em,
    roffset=1em,
    toffset=0.5em,
    boffset=0.5em,
]

% Tipp-Box (grün)
\defineframedtext[tippbox][
    frame=on,
    framecolor=secondarycolor,
    background=color,
    backgroundcolor=secondarycolor:1,
    width=\textwidth,
    before={\blank[medium]},
    after={\blank[medium]},
    loffset=1em,
    roffset=1em,
    toffset=0.5em,
    boffset=0.5em,
]

% Warnung-Box (rot)
\defineframedtext[warnbox][
    frame=on,
    framecolor=warningcolor,
    background=color,
    backgroundcolor=warningcolor:1,
    width=\textwidth,
    before={\blank[medium]},
    after={\blank[medium]},
    loffset=1em,
    roffset=1em,
    toffset=0.5em,
    boffset=0.5em,
]

% =============================================================================
% AUFGABEN UND ÜBUNGEN
% =============================================================================

% Aufgaben-Box (orange)
\defineframedtext[aufgabenbox][
    frame=on,
    framecolor=accentcolor,
    background=color,
    backgroundcolor=accentcolor:1,
    width=\textwidth,
    before={\blank[medium]},
    after={\blank[medium]},
    loffset=1em,
    roffset=1em,
    toffset=0.5em,
    boffset=0.5em,
]

% Übungs-Box (grün)
\defineframedtext[uebungsbox][
    frame=on,
    framecolor=secondarycolor,
    background=color,
    backgroundcolor=secondarycolor:1,
    width=\textwidth,
    before={\blank[medium]},
    after={\blank[medium]},
    loffset=1em,
    roffset=1em,
    toffset=0.5em,
    boffset=0.5em,
]

% =============================================================================
% LERNZIELE
% =============================================================================

% Lernziele-Box (grün)
\defineframedtext[lernziele][
    frame=off,
    background=color,
    backgroundcolor=secondarycolor:1,
    width=\textwidth,
    before={\blank[medium]},
    after={\blank[medium]},
    loffset=1em,
    roffset=1em,
    toffset=0.5em,
    boffset=0.5em,
]

% =============================================================================
% MARKDOWN-ÄHNLICHE SHORTCUTS
% =============================================================================

% Fett
\define[1]\strong{{\bf #1}}

% Kursiv
\define[1]\emph{{\it #1}}

% Durchgestrichen
\define[1]\strike{{\overstrike #1}}

